\subsection{Irreducibility of cyclotomic polynomials}

Let p be the characteristic exponent of K and let n be an integer prime to p.~We denote by \(\Phi_{n} \in \mathbf{K}[\mathbf{X}]\) the image of the polynomial with integer coefficients \(\Phi_{n}\) under the unique homomorphism of Z{[}X{]} into K{[}X{]} which maps X to X.

Lemma 3. --- The roots of \(\Phi_n\) in \(\mathbf{K}_s\) are the primitive \(n\) -th roots of unity.

Denote by \(S_{n}\) the set of roots of \(\Phi_{n}\) in \(K_{s}\) . By Formula (6), the set \(\mu_{n}(K_{s})\) is the union of the \(S_{d}\) for d dividing n.~Every primitive n-th root of unity thus belongs to \(S_{n}\) and the lemma now follows from Prop. 4 (V, p.~81).

PROPOSITION 6. --- Let p be the characteristic exponent of K and let \(n \geqslant 1\) be an integer prime to p.~For the polynomial \(\Phi_{n}(\mathbf{X})\) to be irreducible in K{[}X{]} it is necessary and sufficient that the homomorphism \(\chi_{n}: \operatorname{Gal}(K_{s}/K) \to (\mathbf{Z}/n\mathbf{Z})^{*}\) should be surjective.

By Lemma 3 we have \(\mathbf{R}_n(\mathbf{K}) = \mathbf{K}(\zeta)\) for each root \(\zeta\) of \(\Phi_n(\mathbf{X})\) and hence \(\Phi_n(\mathbf{X})\) is irreducible in \(\mathbf{K}[\mathbf{X}]\) if and only if the degree \(\varphi(n)\) of \(\Phi_n(\mathbf{X})\) is equal to \([\mathbf{R}_n(\mathbf{K}) : \mathbf{K}]\) . Further, the Galois group of \(\mathbf{R}_n(\mathbf{K})\) over \(\mathbf{K}\) is of order \([\mathbf{R}_n(\mathbf{K}) : \mathbf{K}]\) and it is isomorphic to the subgroup of \((\mathbf{Z}/n\mathbf{Z})^*\) which is the image of \(\chi_n\) . Now Prop. 6 follows from the fact that \((\mathbf{Z}/n\mathbf{Z})^*\) is of order \(\varphi(n)\) .

THEOREM 2 (Gauss). --- Let \(\bar{\mathbf{Q}}\) be an algebraic closure of \(\mathbf{Q}\) and let \(n \geqslant 1\) be an integer.

\begin{enumerate}
\def\labelenumi{\alph{enumi})}
\item
  The cyclotomic polynomial \(\Phi_n(\mathbf{X})\) is irreducible in \(\mathbf{Q}[\mathbf{X}]\) .
\item
  The degree of \(\mathbf{R}_n(\mathbf{Q})\) over \(\mathbf{Q}\) is \(\varphi(n)\) .
\item
  The homomorphism \(\chi_{n}\) of \(\operatorname{Gal}(\bar{\mathbf{Q}} / \mathbf{Q})\) into \((\mathbf{Z} / n\mathbf{Z})^{*}\) is surjective and defines by passage to quotients an isomorphism of \(\operatorname{Gal}(\mathrm{R}_n(\mathbf{Q}) / \mathbf{Q})\) onto \((\mathbf{Z} / n\mathbf{Z})^{*}\) .
\end{enumerate}

Taking account of Prop. 6, we need only prove \(c\) ). Every integer \(r\) prime to \(n\) is a product of prime numbers \(p_1, \ldots, p_s\) not dividing \(n\) ; so it is enough to show that for every prime number \(p\) not dividing \(n\) , the mapping \(x \mapsto x^p\) of \(\mu_n(\mathbf{Q})\) into itself extends to an automorphism of \(\mathbf{R}_n(\mathbf{Q})\) . It suffices to prove that if \(\zeta\) is a primitive \(n\) -th root of unity, the minimal polynomial \(f\) of \(\zeta\) over \(\mathbf{Q}\) is equal to the minimal polynomial \(g\) of \(\zeta^p\) over \(\mathbf{Q}\) .

We shall assume that \(f \neq g\) and argue by contradiction. The polynomials \(f\) and \(g\) are monic irreducible in \(\mathbf{Q}[\mathbf{X}]\) and divide \(\mathbf{X}^n - 1\) , so there exists \(u \in \mathbf{Q}[\mathbf{X}]\) such that \(\mathbf{X}^n - 1 = fgu\) (IV, p.~13, Prop. 13). Lemma 2 (V, p.~82) shows that \(f, g\) and \(u\) have integer coefficients. Let us denote by \(\bar{v}\) the polynomial with coefficients in \(\mathbf{F}_p\) obtained from a polynomial \(v \in \mathbf{Z}[\mathbf{X}]\) by reduction mod \(p\) . We thus have \(\mathbf{X}^n - 1 = \bar{f}\bar{g}\bar{u}\) in \(\mathbf{F}_p[\mathbf{X}]\) .

Further we have \(g(\zeta^p) = 0\) and so \(g(\mathbf{X}^p)\) is a multiple of \(f(\mathbf{X})\) in \(\mathbf{Q}[\mathbf{X}]\) . By Lemma 2 there exists \(h \in \mathbf{Z}[\mathbf{X}]\) such that \(g(\mathbf{X}^p) = f(\mathbf{X}) \cdot h(\mathbf{X})\) . Now we have \(v(\mathbf{X}^p) = v(\mathbf{X})^p\) for every polynomial \(v \in \mathbf{F}_p[\mathbf{X}]\) . By reduction mod \(p\) we thus obtain \(\overline{g}^p = \overline{f}\overline{h}\) . If \(v\) is an irreducible polynomial in \(\mathbf{F}_p[\mathbf{X}]\) dividing \(\overline{f}\) , it must then divide \(\overline{g}\) . Since \(\overline{f}\overline{g}\) divides \(\mathbf{X}^n - 1\) , we conclude that \(v^2\) divides \(\mathbf{X}^n - 1\) in \(\mathbf{F}_p[\mathbf{X}]\) . This is absurd because the polynomial \(\mathbf{X}^n - 1\) is separable in \(\mathbf{F}_p[\mathbf{X}]\) .

It may be shown that for every abelian extension E of finite degree over Q there exists an integer \(n \geqslant 1\) such that E is isomorphic to a subextension of \(\mathrm{R}_{n}(\mathbf{Q})\) . * In other words, the field \(\mathbf{Q}(\mu_{\infty}(\mathbf{C}))\) is an abelian closure of \(\mathbf{Q}\) . *(« Kronecker-Weber theorem »).

